\documentclass[10pt,a4paper]{amsart}
\usepackage[margin=19mm]{geometry}
\usepackage{amsmath,amssymb,mathtools}
\usepackage[scr=rsfs,cal=euler]{mathalfa}
\usepackage{xfrac}
\usepackage[unicode,hidelinks,bookmarks=false]{hyperref}
\newcommand{\ie}{i.e.\ }
\newcommand{\cf}{cf.\ }
\newcommand{\resp}{respectively\ }
\newcommand{\Subsection}{\subsection}
\providecommand{\nobreakdash}{\nobreak\hbox{-}}
\setlength{\emergencystretch}{2em}
\multlinegap0pt
\usepackage{bm}
\usepackage{bbm}
\usepackage{dsfont}
\usepackage{xspace}
\usepackage{cancel}
\usepackage{mathtools}
\usepackage[shortlabels]{enumitem}
\usepackage{amsthm}
\makeatletter
\@ifundefined{theorem}{\newtheorem{theorem}{Theorem}}{}
\@ifundefined{corollary}{\newtheorem{corollary}[theorem]{Corollary}}{}
\@ifundefined{proposition}{\newtheorem{proposition}[theorem]{Proposition}}{}
\makeatother
\DeclareMathOperator{\diam}{diam}
\newcommand{\C}{\mathbb C}
\newcommand{\D}{\mathbb D}
\newcommand{\N}{\mathbb N}
\newcommand{\R}{\mathbb R}
\renewcommand{\bar}{\overline}
\title[Characterization of tangent quasicircles and quasiannuli]{Characterization of tangent quasicircles and quasiannuli}
\author{Dimitrios Ntalampekos}
\date{Source: arXiv:2602.17236v1 (CC BY 4.0)}
\begin{document}
\maketitle

\noindent\emph{Source and licence.} Characterization of tangent quasicircles and quasiannuli. Version arXiv:2602.17236v1 (CC BY 4.0), distributed under the Creative Commons Attribution 4.0 International licence, as recorded in the arXiv licence metadata for this exact version and shown on the abstract page https://arxiv.org/abs/2602.17236v1. Full rights notice: licenses/arxiv-2602.17236-CC-BY-4.0.txt.\par\smallskip

\noindent\textit{Editorial scope.} Counted unit: the Proposition whose source label is \texttt{corollary:extension\_annulus} in the article (source0.tex lines 1683--1685, source label \texttt{corollary:extension\_annulus}), delivered together with the article's own complete original proof of it (source0.tex lines 1687--1736, one \texttt{proof} environment of 50 lines). No mathematics is written, repaired or re-derived: statement and proof are the article's own text line for line. Required context retained uncounted: definition:qs (lines 526--528); corollary:beurling\_ahlfors (lines 621--627); theorem:extension\_annulus (lines 1636--1643). Every definition, display and label of those lines is retained unchanged. Omitted from this package and not redistributed: 1--525, 529--620, 628--1635, 1644--1682, 1686--1686, 1737--1847. Nothing inside a retained excerpt is omitted or altered: every retained body is delivered exactly as the article's lines are written, so no omission receipt of any kind accompanies this package. Citations: the retained excerpts carry no citation command at all, so no bibliography entry of the article is transcribed and no reference is redistributed. Reference adaptation: none was needed and none was made; every reference inside the delivered unit resolves inside it, so no unresolved reference or citation is produced. Printed slips kept literal: no printed word, symbol, hypothesis or number of the counted statement or of its proof was altered. Layout and class: the article's own class is \texttt{amsart} with packages including amsrefs, amssymb, mathtools, mathrsfs, enumitem, color, comment, pgfplots, hyperref; the wrapper is the lane's pinned \texttt{amsart} editorial wrapper at 10pt on A4, which restates the theorem-like environments the retained text needs and adds the article's own macro definitions that the retained text needs (\texttt{\textbackslash C}, \texttt{\textbackslash D}, \texttt{\textbackslash N}, \texttt{\textbackslash R}, \texttt{\textbackslash bar}, \texttt{\textbackslash diam}), with extra wrapper packages (bm, bbm, dsfont, xspace, cancel, mathtools, enumitem). The delivered numbering is the unit's own. Assets: no retained span contains a figure environment, a graphics-inclusion command, a PDF-page inclusion, a table or a TikZ picture; no image file is copied into this package, the package declares no asset dependency and no image file is redistributed with it. Licence: version arXiv:2602.17236v1 of this article is distributed under the Creative Commons Attribution 4.0 International licence, as recorded in the arXiv licence metadata for this exact version and shown on the abstract page https://arxiv.org/abs/2602.17236v1. A full rights notice is delivered at licenses/arxiv-2602.17236-CC-BY-4.0.txt. Characters: every retained excerpt is pure ASCII, so the delivered engine, XeLaTeX with its default Latin Modern text font, reports no missing character.
\begin{align}\label{definition:qs}
\frac{1}{2\eta\left( \frac{\diam (B)}{\diam (A)}\right)} \leq \frac{\diam (f(A))}{\diam (f(B))}\leq \eta\left( \frac{2\diam (A)}{\diam (B)}\right).
\end{align}

\begin{corollary}\label{corollary:beurling_ahlfors}
Let $f\colon \partial \D\to \partial \D$ be an orientation-preserving homeomorphism.
\begin{enumerate}[label=\normalfont(\arabic*)]
\item\label{corollary:beurling_ahlfors:1} If $f$ is quasi-M\"obius, then there exists an extension of $f$ to a quasiconformal homeomorphism of $\D$, quantitatively.
\item\label{corollary:beurling_ahlfors:2} If $f$ quasisymmetric, then there exists an extension of $f$ to a quasisymmetric homeomorphism of $\D$ that fixes the point $0$, quantitatively.
\end{enumerate}
\end{corollary}

\begin{proposition}\label{theorem:extension_annulus}
Let $\eta$ be a distortion function and let $N\in \N$ with $N\geq 2$ and $a\geq 1$. Let $h\colon  \mathbb S(0,1)\cup \mathbb S(0,e^{2\pi/N}) \to \mathbb S(0,1)\cup \mathbb S(0,e^{2\pi/N}) $ be a homeomorphism that is orientation-pre\-serving and preserves each of the two circles. Suppose that for each $t\in \R$
\begin{enumerate}[label=\normalfont(\arabic*)]
\item the map $h$ is $\eta$-quasisymmetric in $B_r(t)=\{re^{i\theta}: t \frac{2\pi}{N}\leq \theta\leq t\frac{2\pi}{N}+ \frac{2\pi}{N}\}$ for $r=1$ and $r=e^{2\pi/N}$ and
\item the four endpoints of the arcs $B_1(t), B_{e^{2\pi/N}}(t)$ are mapped to points with mutual distances in the interval $[a^{-1} N^{-1},a N^{-1}]$. 
\end{enumerate}
Then there exists an extension of $h$ to a homeomorphism $H\colon \bar{\mathbb A}(0;1,e^{2\pi/N})\to \bar{\mathbb A}(0;1,e^{2\pi/N})$ that is $K(\eta,a)$-quasiconformal in ${\mathbb A}(0;1,e^{2\pi/N})$.
\end{proposition}

\begin{proposition}\label{corollary:extension_annulus}
Let $\eta$ be a distortion function and let $L,L'>1$ such that $L\in (1,M)$ for some $M>1$. Let $h\colon  \mathbb S(0,1)\cup \mathbb S(0,L) \to \mathbb S(0,1)\cup \mathbb S(0,L') $ be an $\eta$-quasisymmetric homeomorphism that is orien\-tation-preserving and preserves each of the two circles. Then there exists an extension of $h$ a $K(\eta,M)$-quasi\-conformal quasiconformal homeomorphism of $\C$.
\end{proposition}

\begin{proof}
Note that $h$ extends quasiconformally to $\D$ and to $\C\setminus \bar \D(0,L)$ by Corollary \ref{corollary:beurling_ahlfors}. Thus, it suffices to extend it in the annulus $\mathbb A(0;1,L)$. Our goal is to modify the map $h$ in a bi-Lipschitz way so that Proposition \ref{theorem:extension_annulus} can be applied. 

To begin with, we show that $L-1\simeq_{\eta,M} L'-1$. Let $n\in \N$ such that 
$$\frac{M-1}{n+1}\leq  L-1 <  \frac{M-1}{n}.$$
Consider a partition  $\{z_0,\dots,z_n\}$ of $\mathbb S(0,1)$ such that the arc $A_i$ from $z_{i-1}$ to $z_i$ has length $2\pi/(n+1)$ for $i\in \{0,\dots,n\}$, where $z_{-1}\coloneqq z_n$. Thus, for $i\in \{0,\dots,n\}$, we have
$$|z_{i-1}-z_i|\leq \frac{2\pi}{n+1}.$$
Fix $0\in \{1,\dots,n\}$. Denote by $B_i$ be the complementary arc of $A_i$ and note that $\diam_e(B_i)=2=\diam_e(\mathbb S(0,1))$. Thus, $\diam_e(h(B_i)) \geq \eta(1)^{-1}$ by \eqref{definition:qs}. This implies that $\ell(h(A_i))\leq 2\pi -\eta(1)^{-1}$, so
\begin{align*}
\ell(h(A_i))\leq C_1(\eta)|h(z_{i-1})-h(z_i)|.
\end{align*}
Also,
\begin{align*}
\frac{|h(z_{i-1})-h(z_i)|}{L'-1}&=\frac{|h(z_{i-1})-h(z_i)|}{|h(z_{i})- L'h(z_i)|}\\
&\leq \eta\left( \frac{|z_{i-1}-z_i|}{|z_i- h^{-1}(Lh(z_i))|}\right)\leq \eta \left(\frac{ 2\pi(n+1)^{-1}}{L-1}\right)\leq \eta\left(\frac{2\pi}{M-1}\right). 
\end{align*}
Therefore, $\ell(h(A_i))\leq C_2(\eta,M) (L'-1)$. We sum over $i\in \{0,\dots,n\}$ and we obtain
$$2\pi \leq C_2(\eta,M) \cdot (n+1) (L'-1)\leq C_2(\eta,M) \cdot 2n  (L'-1) \leq C_2(\eta,M) 2(M-1) \frac{L'-1}{L-1}.$$
If we show that $L'\lesssim_{\eta,M}1$, then we can apply the same argument to $h^{-1}$ in order to obtain an inequality in the reverse direction. Observe that by applying \eqref{definition:qs} to $A=\mathbb S(0,1)$ and $B=\mathbb S(0,1)\cup \mathbb S(0,L)$, we have
$$\frac{2}{2L'} \geq \frac{1}{2\eta\left(\frac{2L}{2}\right)}> \frac{1}{2\eta(M)}.$$ 
Therefore, $L'< 2\eta(M)\eqqcolon M'$, as desired. This completes the proof of the claim that $L-1\simeq_{\eta,M} L'-1$.

Next, we modify the map $h$ in a bi-Lipschitz way. Let $N\in \N$, $N\geq 2$, such that $\frac{M-1}{N}\leq  L-1<\frac{M-1}{N-1}$, so $L-1\simeq_M\frac{1}{N}$. Note that $\frac{1}{N}\simeq e^{2\pi/N}-1$, so  
$$L-1\simeq_{M} e^{2\pi/N}-1\quad \text{and}\quad L'-1\simeq_{\eta,M}e^{2\pi/N}-1.$$
Consider the maps
\begin{align*}
\sigma(re^{i\theta}) &= \bigg(1+ \frac{e^{2\pi/N}-1}{L-1}(r-1)\bigg) e^{i\theta}\,\,\, \text{and} \,\,\, \tau(re^{i\theta}) &= \bigg(1+ \frac{e^{2\pi/N}-1}{L'-1}(r-1) \bigg) e^{i\theta}.
\end{align*}
The map $\sigma$ (resp.\ $\tau$) interpolates between the identity map on $\mathbb S(0,1)$ and a scaling on $\mathbb S(0,L)$ (resp.\ $\mathbb S(0,L')$). The map $\sigma$ is  bi-Lipschitz on $\bar{ \mathbb A}(0;1,L)$, depending only on $M$, and maps it onto $\bar{\mathbb A}(0;1,e^{2\pi/N})$. The map $\tau$ is bi-Lipschitz on $\bar{ \mathbb A}(0;1,L')$, depending only on $\eta,M$, and maps it onto $\bar{\mathbb A}(0;1,e^{2\pi/N})$. If we can extend the map $\tau \circ h\circ \sigma^{-1}$ to a $K(\eta,M)$-quasiconformal map of the annulus $\mathbb A(0;1,e^{2\pi/N})$, then we can also extend the original map $h$ to a $K'(\eta,M)$-quasiconformal map from $\mathbb A(0;1,L)$ onto $\mathbb A(0;1,L')$. Therefore, by replacing the map $h$ with $\tau \circ h\circ \sigma^{-1}$, we may assume that $L=L'=e^{2\pi/N}\in (1,e^{\pi}]$ and $h$ is $\eta$-quasisymmetric.

We will verify the assumptions of Proposition \ref{theorem:extension_annulus}. Consider the curvilinear square $S(t)=\{z\in \bar {\mathbb A}(0;1,L): t\log L\leq   \arg(z)\leq (t+1)\log L\}$. It suffices to show that its vertices are mapped to points with mutual distances in the interval $[a^{-1}\log L, a\log L]$ for some $a=a(\eta)\geq 1$.

Let $z,w\in \mathbb S(0,1)$ be two vertices of $S(t)$. The other two vertices are $Lz$ and $Lw$. We will bound the distances $|h(z)-h(w)|$, $|h(z)-h(Lz)|$, $|h(z)-h(Lw)|$, $|h(Lz)-h(Lw)|$. Then the bounds for $|h(w)-h(Lw)|$ and $|h(w)-h(Lz)|$ follow by symmetry. 

Note that the arc of $S(t)$ between $z$ and $w$ has length $\log L\leq \pi$. Thus, $ \frac{2}{\pi}\log L \leq |z-w|\leq \log L$. Moreover, since $L\in (1,e^{\pi}]$, we have $c_0^{-1}(L-1)\leq \log L \leq L-1$ for $c_0=\pi^{-1}(e^\pi-1)$. Therefore, we have
$$\frac{L-1}{|h(z)-h(w)|}\leq \frac{|h(Lz)-h(z)|}{|h(z)-h(w)|}\leq \eta\left( \frac{|Lz-z|}{|z-w|}\right) \leq \eta\left(\frac{L-1}{2\pi^{-1}\log L}\right)\leq \eta\left(\frac{c_0\pi}{2}\right) $$ 
and
\begin{align*}
\frac{|h(z)-h(w)|}{L-1}&=\frac{|h(z)-h(w)|}{|h(z)- Lh(z)|}\leq \eta\left( \frac{|z-w|}{|z- h^{-1}(Lh(z))|}\right)\leq \eta \left(\frac{\log L}{L-1}\right)\leq \eta(1). 
\end{align*}
Thus $L-1\leq |h(z)-h(Lz)|\leq \eta(\frac{c_0\pi}{2})\eta(1)(L-1)$ and $\eta(\frac{c_0\pi}{2})^{-1}(L-1)\leq  |h(z)-h(w)|\leq \eta(1)(L-1)$.  These bounds give immediately bounds for $|h(z)-h(Lw)|$.  Similarly, 
\begin{align*}
\frac{|h(Lz)-h(Lw)|}{L-1}& = \frac{|h(Lz)-h(Lw)|}{|h(z)-Lh(z)|} \leq \eta \left(\frac{L|z-w|}{L-1}\right)\leq \eta(L)\leq \eta(e^\pi)
\end{align*}
and
\begin{align*}
\frac{L-1}{|h(Lz)-h(Lw)|}\leq \frac{|h(Lz)-h(z)|}{|h(Lz)-h(Lw)|} \leq \eta\left(\frac{L-1}{2\pi^{-1}L\log L}\right)\leq \eta\left(\frac{c_0\pi}{2L}\right)\leq \eta\left(\frac{c_0\pi}{2}\right).
\end{align*}
So, $\eta(\frac{c_0\pi}{2})^{-1}(L-1)\leq |h(Lz)-h(Lw)|\leq \eta(e^\pi)(L-1)$. This completes the proof.
\end{proof}

\end{document}
